% SPDX-License-Identifier: Apache-2.0
\section{A Function on a Literal}
\label{sec:x-litfn}

A function under \code{\#[lower]} may be called on a literal as well
as on a signal, and then what it reads of its argument is part of a
constant. An SD card checks every command with a CRC-7, shifted in a
bit at a time, and the first bit of a command goes into an empty
register, so the natural way to write it is the step applied to a
zero: \code{crc7\_step(U::<7>::from(0u8), b)}. The step reads the
register's top bit, and inlined on a literal that is a bit of the
literal.

VHDL does not index a qualified expression. The lowering used to
write \code{unsigned'("0000000")(6)}, which Verilator analysed and
nvc refused, so the netlist was right in one language only; that is
issue 555. A bit or a part of a literal at a known place is now
folded into the literal bit or bits it names where the expression is
built, so neither netlist indexes one: the first bit's step reads
\code{1'b0}, and \code{top2} of \code{0x60} is \code{2'b11}.

The run shifts in SD's \code{CMD0}, forty bits, twice, so that the
second command's first bit goes through the step on the literal while
the register still holds what the first command left. Each time the CRC
is \code{0x4A}, which is what the SD specification gives for
\code{CMD0}, and the run asserts it. The netlist is checked against
the run under both simulators.

\begin{figure}[htbp]
\centering
\input{litfn_timing.tex}
\caption{The end of the first \code{CMD0} and the start of the second:
\code{sum} reaches \code{0x4A} as the fortieth bit goes in, and
\code{start} then marks the second command's first bit, which goes
through the step on a literal zero while the register still holds
what the first command left, and the register starts again from zero.}
\label{fig:x-litfn}
\end{figure}

\lstinputlisting[caption={\code{ex\_litfn.rs}. Run with \code{bazel run //lib/examples:ex\_litfn}.},label={lst:x-litfn}]{lib/examples/ex_litfn.rs}

\lstinputlisting[style=ascii,caption={What \code{ex\_litfn} printed when the build ran it: the CRC of each round, then the lowering, the literal's bits folded.},label={lst:x-litfn-out}]{out_litfn.txt}
